%=====================================================================
\chapter{Supervised Learning II: Classification}\label{ch:classification}
%=====================================================================
\locator{3}{Part III --- Machine Learning Core}

\begin{outcomes}
By the end of this chapter you will be able to:
\begin{tight}
  \item \textbf{Explain} how logistic regression converts a linear score into a probability, and interpret its coefficients as odds ratios.
  \item \textbf{Describe} the mechanism, strengths and weaknesses of k-NN, naive Bayes, decision trees and support vector machines.
  \item \textbf{Match} a classifier to a problem's constraints --- data size, interpretability, training cost, feature types.
  \item \textbf{Interpret} a decision boundary plot and relate its shape to model bias.
  \item \textbf{Explain} why the default 0.5 threshold is a choice, not a law.
\end{tight}
\end{outcomes}

\begin{prereq}
\chref{probability} (Bayes' theorem, for naive Bayes), \chref{math} (distance, for
k-NN) and \chref{regression} (the linear score that logistic regression squashes).
\end{prereq}

\begin{keyterms}
classification $\bullet$ decision boundary $\bullet$ logistic (sigmoid) function
$\bullet$ odds ratio $\bullet$ log loss $\bullet$ k-nearest neighbours $\bullet$
naive Bayes $\bullet$ decision tree $\bullet$ Gini impurity $\bullet$ entropy
$\bullet$ support vector machine $\bullet$ margin $\bullet$ kernel $\bullet$
decision threshold
\end{keyterms}

\section{From Regression to Classification}

Classification predicts a \emph{category}, not a number. The natural first attempt
--- fit a straight line and threshold it --- fails, because a line is unbounded and
predicts probabilities above 1 and below 0. Logistic regression fixes this with one
function.

\begin{definitionbox}[Logistic Regression]
Compute a linear score $z = w_0 + \sum_j w_j x_j$, then squash it into $(0,1)$
with the \term{sigmoid}:
\[
\hat{p} \;=\; \sigma(z) \;=\; \frac{1}{1 + e^{-z}}
\]
Predict the positive class when $\hat{p}$ exceeds a chosen threshold. Despite its
name it is a classifier, not a regressor.
\end{definitionbox}

\dsfig{%
\begin{tikzpicture}
\begin{groupplot}[
  group style={group size=2 by 1, horizontal sep=1.5cm},
  width=6.6cm, height=4.6cm,
  axis lines=left, tick label style={font=\tiny}, label style={font=\tiny},
  title style={font=\scriptsize\bfseries},
]
\nextgroupplot[title={\color{accent}A line cannot be a probability},
  xlabel={engagement score}, ylabel={outcome / prediction},
  xmin=0, xmax=10, ymin=-0.45, ymax=1.45]
\addplot[only marks, mark=*, mark size=1.3pt, primary] coordinates {
 (0.8,0)(1.5,0)(2.2,0)(2.9,0)(3.6,0)(4.3,1)(5.0,0)(5.7,1)(6.4,1)(7.1,1)(7.8,1)(8.5,1)(9.2,1)};
\addplot[accent, line width=1.2pt, domain=0:10] {-0.22 + 0.135*x};
\draw[gray!55, dashed] (axis cs:0,1) -- (axis cs:10,1);
\draw[gray!55, dashed] (axis cs:0,0) -- (axis cs:10,0);
\node[font=\tiny, accent!80!black, anchor=west] at (axis cs:5.0,1.3) {predicts $>1$};
\node[font=\tiny, accent!80!black, anchor=west] at (axis cs:0.2,-0.32) {predicts $<0$};

\nextgroupplot[title={\color{greenhl}The sigmoid always lies in $(0,1)$},
  xlabel={linear score $z$}, ylabel={$\sigma(z)$},
  xmin=-6, xmax=6, ymin=-0.08, ymax=1.12]
\addplot[greenhl, line width=1.4pt, domain=-6:6, samples=120] {1/(1+exp(-x))};
\draw[gray!55, dashed] (axis cs:-6,0.5) -- (axis cs:6,0.5);
\draw[gray!55, dashed] (axis cs:0,-0.08) -- (axis cs:0,0.5);
\addplot[accent, only marks, mark=*, mark size=1.8pt] coordinates {(0,0.5)};
\node[font=\tiny, accent, anchor=west] at (axis cs:0.25,0.42) {$z=0 \Rightarrow p=0.5$};
\node[font=\tiny, text=gray!55!black, anchor=west] at (axis cs:-5.8,0.93) {saturates at 1};
\node[font=\tiny, text=gray!55!black, anchor=west] at (axis cs:1.6,0.07) {saturates at 0};
\end{groupplot}
\end{tikzpicture}}%
{Why logistic regression exists. A linear fit to a 0/1 outcome produces impossible
probabilities; the sigmoid maps any real score into a valid probability and saturates
gracefully at both ends.}{sigmoid}

\begin{worked}[Interpreting a Logistic Coefficient]
A model predicting course failure returns
\[
z = -1.2 - 0.045\,(\text{attendance\%}) + 0.62\,(\text{days\_absent})
\]

\smallskip
\textbf{Step 1 --- a concrete prediction.} For a student with 60\% attendance and
4 days absent:
\[
z = -1.2 - 0.045(60) + 0.62(4) = -1.2 - 2.7 + 2.48 = -1.42
\]
\[
\hat{p} = \frac{1}{1+e^{1.42}} = \frac{1}{1+4.137} = 0.195
\]
About a 19.5\% estimated probability of failing.

\smallskip
\textbf{Step 2 --- the odds ratio.} Exponentiating a coefficient gives the
multiplicative effect on the \emph{odds}:
\[
e^{0.62} = 1.86
\]
Each additional day absent multiplies the odds of failing by 1.86 --- an 86\%
increase in odds, holding attendance constant. This is how logistic models are
reported in practice, because ``the odds nearly double per day absent'' is
comprehensible in a way that ``the log-odds rise by 0.62'' is not.

\smallskip
\textbf{Step 3 --- the trap.} Odds are not probability. Odds rising 86\% moves a
probability of 0.05 to 0.089, but a probability of 0.5 to 0.65. The same
coefficient has a very different effect depending on where the student started.
\end{worked}

\section{Four More Classifiers}

\dsfig{%
\begin{tikzpicture}[font=\scriptsize,
  hd/.style={rounded corners=3pt, fill=#1, text=white, font=\scriptsize\bfseries,
             text width=4.35cm, align=center, minimum height=0.62cm},
  bd/.style={rounded corners=3pt, draw=#1, fill=#1!7, text=#1!50!black,
             text width=4.35cm, align=left, font=\scriptsize, minimum height=2.35cm,
             inner sep=5pt}]

\foreach \i/\c/\name/\body [count=\n from 0] in {%
  0/secondary/{Logistic Regression}/{\textbf{Idea:} squash a weighted sum.\\[3pt]
     \textbf{Good:} fast, gives calibrated probabilities, coefficients explain
     themselves.\\[3pt] \textbf{Bad:} boundary must be a straight line.},
  1/greenhl/{k-Nearest Neighbours}/{\textbf{Idea:} vote among the $k$ closest
     points.\\[3pt] \textbf{Good:} no training at all; any boundary shape.\\[3pt]
     \textbf{Bad:} slow at prediction, needs scaling, dies in high dimensions.},
  2/goldhl/{Naive Bayes}/{\textbf{Idea:} Bayes' theorem, pretending features are
     independent.\\[3pt] \textbf{Good:} tiny data, very fast, excellent for
     text.\\[3pt] \textbf{Bad:} the independence assumption is usually false.},
  3/purplehl/{Decision Tree}/{\textbf{Idea:} ask a sequence of yes/no questions.\\[3pt]
     \textbf{Good:} readable by non-experts, no scaling, mixed data types.\\[3pt]
     \textbf{Bad:} overfits badly alone --- see \chref{features}.},
  4/accent/{Support Vector Machine}/{\textbf{Idea:} the boundary with the widest
     margin.\\[3pt] \textbf{Good:} strong with many features, kernels give
     curves.\\[3pt] \textbf{Bad:} slow on large $n$, no natural probabilities.}}{%
  % 3 across on the top row, 2 on the second -- 4.35cm cards stay legible
  \pgfmathsetmacro{\xx}{mod(\n,3)*4.7}
  \pgfmathsetmacro{\yy}{-floor(\n/3)*3.55}
  \node[hd=\c] at (\xx,\yy) {\name};
  \node[bd=\c] at (\xx,\yy-1.6) {\body};}
\end{tikzpicture}}%
{The five classical classifiers, with the trade-off that decides between them. There is no
best one --- \chref{mlfundamentals}'s no-free-lunch result in practical form.}{classifier-cards}

\subsection{Decision Trees: How the Split Is Chosen}

\begin{definitionbox}[Gini Impurity]
For a node containing class proportions $p_k$,
$\;G = 1 - \sum_k p_k^2$. It is 0 when the node is pure (one class only) and
maximal when classes are evenly mixed. A tree chooses the split that reduces
weighted impurity the most.
\end{definitionbox}

\begin{worked}[Choosing a Split by Gini]
100 students, 30 of whom failed. Root impurity:
\[
G_{\text{root}} = 1 - (0.7^2 + 0.3^2) = 1 - (0.49 + 0.09) = 0.42
\]

\textbf{Candidate split: attendance $< 50\%$?}
\begin{tight}
  \item Left (attendance $<50$): 25 students, 20 failed.
        $G_L = 1 - (0.2^2 + 0.8^2) = 1 - (0.04+0.64) = 0.32$
  \item Right: 75 students, 10 failed.
        $G_R = 1 - (0.867^2 + 0.133^2) = 1 - (0.751+0.018) = 0.231$
\end{tight}
Weighted impurity: $\frac{25}{100}(0.32) + \frac{75}{100}(0.231) = 0.08 + 0.173 = 0.253$.

\textbf{Gain:} $0.42 - 0.253 = 0.167$.

\smallskip
\textbf{Candidate split: posted on forum?}
Left: 40 students, 14 failed, so $G_L = 1-(0.35^2+0.65^2) = 0.455$.
Right: 60 students, 16 failed, so $G_R = 1-(0.267^2+0.733^2) = 0.391$.
Weighted impurity $\frac{40}{100}(0.455) + \frac{60}{100}(0.391) = 0.417$,
a gain of only $0.42 - 0.417 = 0.003$.

\smallskip
\textbf{The tree picks attendance}, because it separates the classes far better.
Repeat recursively on each child, and that is the entire algorithm.
\end{worked}

\section{Decision Boundaries}

The single most useful way to understand a classifier is to see the shape of
boundary it is capable of drawing.

\dsfig{%
\begin{tikzpicture}
\begin{groupplot}[
  group style={group size=4 by 1, horizontal sep=0.55cm},
  width=4.3cm, height=3.9cm,
  axis lines=left, xtick=\empty, ytick=\empty,
  title style={font=\tiny\bfseries},
  xmin=0, xmax=10, ymin=0, ymax=10,
]
% shared point sets: two interleaved clusters
\nextgroupplot[title={\color{secondary}Logistic: a line}]
\addplot[secondary!25, fill=secondary!12, draw=none] coordinates {(0,0)(10,7.6)(10,0)} \closedcycle;
\addplot[only marks, mark=*, mark size=1pt, primary] coordinates {
 (1.2,1.4)(2.1,2.6)(1.8,3.9)(3.2,1.8)(2.6,4.6)(4.1,2.9)(3.6,5.4)(1.4,5.8)(4.8,4.2)(2.9,6.7)};
\addplot[only marks, mark=triangle*, mark size=1.3pt, accent] coordinates {
 (7.2,8.4)(8.1,6.6)(6.8,5.9)(8.6,8.8)(5.9,7.8)(9.1,7.2)(7.6,4.9)(6.2,9.1)(8.9,5.4)(5.4,6.2)};
\addplot[secondary, line width=1.1pt] coordinates {(0,0)(10,7.6)};

\nextgroupplot[title={\color{greenhl}k-NN: wiggly}]
\addplot[only marks, mark=*, mark size=1pt, primary] coordinates {
 (1.2,1.4)(2.1,2.6)(1.8,3.9)(3.2,1.8)(2.6,4.6)(4.1,2.9)(3.6,5.4)(1.4,5.8)(4.8,4.2)(2.9,6.7)};
\addplot[only marks, mark=triangle*, mark size=1.3pt, accent] coordinates {
 (7.2,8.4)(8.1,6.6)(6.8,5.9)(8.6,8.8)(5.9,7.8)(9.1,7.2)(7.6,4.9)(6.2,9.1)(8.9,5.4)(5.4,6.2)};
\addplot[greenhl, line width=1.1pt, smooth] coordinates {
 (0,1.2)(1.4,2.4)(2.2,2.0)(3.4,3.6)(4.2,3.2)(5.0,5.0)(5.8,4.4)(6.6,6.0)(7.6,5.6)(8.6,7.2)(10,7.0)};

\nextgroupplot[title={\color{purplehl}Tree: axis-aligned}]
\addplot[only marks, mark=*, mark size=1pt, primary] coordinates {
 (1.2,1.4)(2.1,2.6)(1.8,3.9)(3.2,1.8)(2.6,4.6)(4.1,2.9)(3.6,5.4)(1.4,5.8)(4.8,4.2)(2.9,6.7)};
\addplot[only marks, mark=triangle*, mark size=1.3pt, accent] coordinates {
 (7.2,8.4)(8.1,6.6)(6.8,5.9)(8.6,8.8)(5.9,7.8)(9.1,7.2)(7.6,4.9)(6.2,9.1)(8.9,5.4)(5.4,6.2)};
\addplot[purplehl, line width=1.1pt] coordinates {(5.1,0)(5.1,4.6)};
\addplot[purplehl, line width=1.1pt] coordinates {(5.1,4.6)(10,4.6)};

\nextgroupplot[title={\color{accent}SVM (RBF): smooth curve}]
\addplot[only marks, mark=*, mark size=1pt, primary] coordinates {
 (1.2,1.4)(2.1,2.6)(1.8,3.9)(3.2,1.8)(2.6,4.6)(4.1,2.9)(3.6,5.4)(1.4,5.8)(4.8,4.2)(2.9,6.7)};
\addplot[only marks, mark=triangle*, mark size=1.3pt, accent] coordinates {
 (7.2,8.4)(8.1,6.6)(6.8,5.9)(8.6,8.8)(5.9,7.8)(9.1,7.2)(7.6,4.9)(6.2,9.1)(8.9,5.4)(5.4,6.2)};
\addplot[accent, line width=1.1pt, domain=0:10, samples=60] {1.1 + 0.62*x + 0.9*sin(deg(x/1.7))};
\end{groupplot}
\end{tikzpicture}}%
{The same data, four boundary shapes. Logistic regression can only draw a line; a tree can
only draw axis-parallel steps; k-NN and kernel SVMs can draw anything --- which is powerful
and is also how they overfit.}{decision-boundaries}

\section{The Threshold Is a Decision}

\begin{alertbox}[0.5 Is Not Sacred]
Predicting ``positive when $\hat{p} > 0.5$'' is a default, not a requirement.
Lower the threshold and you catch more true positives at the cost of more false
alarms; raise it and the reverse. On imbalanced problems the default is almost
always wrong. \chref{evaluation} shows how to choose the threshold from the cost
of each error type --- and that is where this chapter's models actually become
usable.
\end{alertbox}

\rowcolors{2}{white}{lightbg}
\begin{longtable}{@{}L{2.7cm}L{2.4cm}L{2.4cm}L{2.5cm}L{4.0cm}@{}}
\toprule
\rowcolor{primary!15}
\textbf{Model} & \textbf{Train speed} & \textbf{Predict speed} & \textbf{Explainable?} & \textbf{Reach for it when} \\
\midrule
\endhead
Logistic regression & Fast & Very fast & \textbf{Yes} --- odds ratios & You must justify decisions to people \\
k-NN & Instant (none) & \textbf{Slow} & Partly & Small data, odd boundary shapes \\
Naive Bayes & Very fast & Very fast & Partly & Text, tiny data, a strong baseline \\
Decision tree & Fast & Very fast & \textbf{Yes} --- readable rules & Mixed types, rules must be auditable \\
SVM (RBF) & Slow on big $n$ & Moderate & No & Many features, moderate $n$, accuracy over clarity \\
\bottomrule
\end{longtable}

\begin{pitfall}
\begin{tight}
  \item \textbf{Using k-NN without scaling.} The largest-range feature becomes the
        only feature (\chref{math}).
  \item \textbf{Believing naive Bayes' probabilities.} Its ranking is often good;
        its probabilities are usually badly calibrated because the independence
        assumption is false.
  \item \textbf{Deploying a single unpruned tree.} It will memorise the training
        set. Use depth limits, or the ensembles of \chref{features}.
  \item \textbf{Reading logistic coefficients as probabilities.} They are
        log-odds; exponentiate them for odds ratios.
\end{tight}
\end{pitfall}

%---------------------------------------------------------------------
\begin{fourlenses}{Choosing a Classifier}
\lensLA{Use \emph{logistic regression or a shallow tree}, even at some cost in
accuracy. A model that flags a student for intervention must be explainable to that
student and to the appeals process --- ``you were flagged because attendance below
50\% and two missed labs'' is defensible; ``the SVM said so'' is not. This is a
requirement from \chref{ethics}, not a preference.}

\lensPM{With 40 sprints of history, \emph{naive Bayes or a depth-2 tree} is the
honest choice. Anything more flexible will fit noise. A two-rule tree that a scrum
master can read out in a stand-up will actually change behaviour, which is the
point --- a marginally better black box that nobody trusts changes nothing.}

\lensIOT{At the edge, prediction cost is a hard constraint: a device with a
microcontroller cannot run k-NN against a million stored points. \emph{Logistic
regression or a small tree} fits in kilobytes and predicts in microseconds. Train
something heavy in the cloud, then distil or replace it with something small enough
to deploy (\chref{cloudedge}).}

\lensSEC{\emph{Naive Bayes} remains a strong, fast baseline for phishing and spam
because text features are numerous and the speed matters at mail-gateway volumes.
For intrusion detection, tree ensembles dominate --- but note that an interpretable
model has a specific operational value here: an analyst must be able to explain to
an incident review why a session was blocked. Also, a k-NN boundary that hugs
training data is trivially evaded by an adversary who probes it.}
\end{fourlenses}

%---------------------------------------------------------------------
\begin{lab}[Five Classifiers, One Comparison]
\begin{lstlisting}[language=Python]
import numpy as np
from sklearn.linear_model import LogisticRegression
from sklearn.neighbors import KNeighborsClassifier
from sklearn.naive_bayes import GaussianNB
from sklearn.tree import DecisionTreeClassifier, export_text
from sklearn.svm import SVC
from sklearn.pipeline import make_pipeline
from sklearn.preprocessing import StandardScaler
from sklearn.model_selection import cross_val_score

models = {
    # scaling matters for the distance/margin-based models
    "logistic":   make_pipeline(StandardScaler(), LogisticRegression(max_iter=1000)),
    "k-NN (k=5)": make_pipeline(StandardScaler(), KNeighborsClassifier(5)),
    "naive Bayes": GaussianNB(),
    "tree (d=3)": DecisionTreeClassifier(max_depth=3, random_state=0),
    "SVM (RBF)":  make_pipeline(StandardScaler(), SVC()),
}

for name, m in models.items():
    s = cross_val_score(m, X_tr, y_tr, cv=5, scoring="roc_auc")
    print(f"{name:12s} AUC {s.mean():.3f} +/- {s.std():.3f}")

# --- the tree's great advantage: you can read it out loud -------------
tree = DecisionTreeClassifier(max_depth=3).fit(X_tr, y_tr)
print(export_text(tree, feature_names=list(feature_names)))

# --- logistic coefficients as odds ratios -----------------------------
lr = make_pipeline(StandardScaler(), LogisticRegression(max_iter=1000)).fit(X_tr, y_tr)
for f, c in zip(feature_names, lr[-1].coef_[0]):
    print(f"  {f:22s} log-odds {c:+.3f}   odds ratio {np.exp(c):.2f}")
\end{lstlisting}
\end{lab}

%---------------------------------------------------------------------
\begin{checkpoint}
\begin{tightnum}
  \item A logistic model has coefficient $+1.1$ on a feature. What is the odds
        ratio, and what does it mean in words?
  \item Why does k-NN require feature scaling while a decision tree does not?
  \item A node contains 40 samples: 40 positive, 0 negative. What is its Gini
        impurity, and what will the tree do next?
\end{tightnum}
\end{checkpoint}

%---------------------------------------------------------------------
\begin{chaptersummary}
\begin{tight}
  \item Logistic regression squashes a linear score through a sigmoid to produce a
        probability; $e^{w_j}$ is the odds ratio for feature $j$.
  \item k-NN needs no training but is slow, scale-sensitive and degrades in high
        dimensions.
  \item Naive Bayes assumes feature independence --- usually false, yet often
        effective, especially on text.
  \item Trees split to reduce Gini impurity, are readable, need no scaling, and
        overfit badly if unconstrained.
  \item SVMs maximise the margin and can draw curved boundaries via kernels, at
        the cost of interpretability and training speed.
  \item Boundary shape is the key intuition: linear, axis-parallel, or arbitrary.
  \item The 0.5 threshold is a default to be revisited in \chref{evaluation}.
\end{tight}
\end{chaptersummary}

\begin{reviewq}
  \item \textbf{(MCQ)} The sigmoid function maps a real number into:
        (a)~$[-1,1]$ (b)~$(0,1)$ (c)~$[0,\infty)$ (d)~$\{0,1\}$.
  \item \textbf{(MCQ)} Which classifier produces a boundary made only of
        axis-parallel segments? (a)~Logistic regression (b)~k-NN (c)~Decision tree
        (d)~RBF SVM.
  \item \textbf{(Short)} Explain what ``naive'' refers to in naive Bayes and give an
        example where the assumption clearly fails.
  \item \textbf{(Short)} A logistic coefficient is $-0.35$. Compute the odds ratio
        and interpret it.
  \item \textbf{(Short)} Why is a single deep decision tree a poor production model,
        and what does \chref{features} do about it?
  \item \textbf{(Applied)} Compute the Gini impurity of a node with 60 samples, 45
        of class A and 15 of class B. Then compute the weighted impurity of a split
        producing children of (30: 28A/2B) and (30: 17A/13B), and state the gain.
  \item \textbf{(Applied)} For each of the four running domains, choose a classifier
        and defend the choice in two sentences, referring to at least one constraint
        that is not accuracy.
\end{reviewq}
